\chapter{La ciudad y direcciones}
\section{Objetivos}
\begin{itemize}[leftmargin=*]
	\item Comprender y dar indicaciones básicas en la ciudad.
	\item Manejar vocabulario de transporte y lugares públicos.
\end{itemize}

\section{Contenidos clave}
\begin{itemize}[leftmargin=*]
	\item Lugares: estación, parada, supermercado, oficina de correos, plaza, calle.
	\item Direcciones: \textit{rechtdoor, linksaf, rechtsaf, tegenover, naast, vlakbij, ver}.
	\item Preguntas modelo: \textit{Waar is ...? Hoe kom ik bij ...? Is het ver?}.
\end{itemize}

\section{Práctica sugerida}
\begin{itemize}[leftmargin=*]
	\item Mapas sencillos: marcar rutas y describirlas en V2.
	\item Diálogo de pedir indicaciones y confirmación.
	\item Señalética: leer y escribir carteles cortos.
\end{itemize}

\section{Microevaluación}
\begin{itemize}[leftmargin=*]
	\item Describe dos rutas distintas de 4 pasos cada una.
\end{itemize}
